\section*{Capítulo \ref{ch:b3:lithosphere-thermal-models} --- Modelos térmicos de la litosfera}

\begin{solution}{exo:b3:lithosphere-thermal-models:1}
$L^2/\kappa = \qty{e12}{s}$, \qty{e14}{s} y \qty{e16}{s}: 0.032, 3.2 y
\qty{317}{Myr}. Un kilómetro olvida su temperatura en decenas de miles de
años; cien kilómetros, en cientos de millones.
\end{solution}

\begin{solution}{exo:b3:lithosphere-thermal-models:2}
$u = \qty{e-3}{m}/\qty{3.156e7}{s} = \qty{3.17e-11}{m.s^{-1}}$, por tanto
$\mathrm{Pe} = 3.17\times10^{-11}\times\num{3e4}/10^{-6} = 0.95$. La advección y
la conducción tienen la misma importancia: la roca ascendente transporta calor hacia arriba
y hace más empinada la geoterma cerca de la superficie, y la conducción por sí sola
subestimaría el gradiente.
\end{solution}

\begin{solution}{exo:b3:lithosphere-thermal-models:3}
$\sqrt{\kappa t} = \sqrt{\num{0.8e-6}\times5\times\num{3.156e13}} =
\qty{11236}{m}$, $\eta = 10000/(2\times11236) = 0.445$, y $T = 1450\times0.471
= \qty{683}{\celsius}$.
\end{solution}

\begin{solution}{exo:b3:lithosphere-thermal-models:4}
$490/\sqrt1 = 490$, $490/5 = 98$ y $490/10 = \qty{49}{mW.m^{-2}}$. La sonda
probablemente subestimaría la primera: la corteza joven está fracturada y abierta a
la circulación de agua marina, que evacua parte del calor a través de manantiales termales.
\end{solution}

\begin{solution}{exo:b3:lithosphere-thermal-models:5}
$z = 2\times1.163\sqrt{\kappa t}$: con $\sqrt{\kappa t} = 15.89$, 31.78 y
\qty{50.25}{km}, $z = 37.0$, 73.9 y \qty{116.9}{km}. La capa enfriada se duplica
de 10 a \qty{40}{Myr}; a los \qty{100}{Myr} ya es más profunda que una
placa de \qty{95}{km}.
\end{solution}

\begin{solution}{exo:b3:lithosphere-thermal-models:6}
$C$ es proporcional a $\alpha$: $\alpha = \num{3.1e-5}\times350/370.5 =
\qty{2.93e-5}{K^{-1}}$. $C$ también es proporcional a $T_m - T_0$ (un manto a
$1450\times350/370.5 = \qty{1370}{\celsius}$) y a $\sqrt\kappa$ (una
difusividad $0.8\times(350/370.5)^2 = \qty{0.71e-6}{m^2.s^{-1}}$): la profundidad por sí sola
no permite separarlos, por lo que también se emplea el flujo de calor.
\end{solution}

\begin{solution}{exo:b3:lithosphere-thermal-models:7}
$3.1\times1450/\num{95e3} = \qty{47.3}{mW.m^{-2}}$ y $3.1\times1350/\num{125e3}
= \qty{33.5}{mW.m^{-2}}$. El semiespacio alcanza esos valores en $(490/47.3)^2 =
\qty{107}{Myr}$ y $(490/33.5)^2 = \qty{214}{Myr}$: la placa delgada y caliente mantiene
el flujo de calor del fondo oceánico antiguo a \qty{47}{mW.m^{-2}} desde aproximadamente \qty{100}{Myr}
en adelante, mientras que la más gruesa y fría permitiría que descendiera por debajo de
\qty{40}{mW.m^{-2}}, valor del semiespacio a los \qty{150}{Myr}.
\end{solution}

\begin{solution}{exo:b3:lithosphere-thermal-models:8}
$A_0h_r = 0.4\times60 = \qty{24}{mW.m^{-2}}$, por tanto $A_0 =
\num{24e-3}/\num{e4} = \qty{2.4}{\micro W.m^{-3}}$ y $q_m = 60 - 24 =
\qty{36}{mW.m^{-2}}$. A \qty{35}{km}, $T = \num{36e-3}\times\num{35e3}/2.5 +
\num{2.4e-6}\times\num{e8}/2.5\times(1 - e^{-3.5}) = 504 + 96\times0.970 =
\qty{597}{\celsius}$.
\end{solution}

\begin{solution}{exo:b3:lithosphere-thermal-models:9}
$h_r = \Delta q_0/\Delta A_0 = \num{18e-3}/\num{1.8e-6} = \qty{10}{km}$, y
$q_m = q_0 - h_rA_0 = 66 - 30 = \qty{36}{mW.m^{-2}}$, ya que $h_rA_0 =
\num{e4}\times\qty{3.0e-6}{W.m^{-3}} = \qty{30}{mW.m^{-2}}$.
\end{solution}

\begin{solution}{exo:b3:lithosphere-thermal-models:10}
(a) $Q(t) = \int_0^tk(T_m - T_0)/\sqrt{\pi\kappa t'}\,\mathrm dt' =
2k(T_m - T_0)\sqrt{t/(\pi\kappa)}$. (b) El calor perdido es el contenido de calor
eliminado de la columna, $Q = \rho_mc\int_0^\infty(T_m - T)\,\mathrm dz$, y
la isostasia da $(\rho_m - \rho_w)(d - d_r) = \rho_m\alpha\int_0^\infty(T_m -
T)\,\mathrm dz = \alpha Q/c$. (c) $c = 3.1/(3300\times\num{0.8e-6}) =
\qty{1174}{J.kg^{-1}.K^{-1}}$; a los \qty{40}{Myr}, $t = \qty{1.262e15}{s}$ y
$Q = 2\times3.1\times1450\times\sqrt{\num{1.262e15}/(\pi\times\num{0.8e-6})} =
\qty{2.015e14}{J.m^{-2}}$; $d - d_r = \num{3.1e-5}\times\num{2.015e14}/
(1174\times2270) = \qty{2343}{m}$, lo que equivale a $370.5\sqrt{40}$.
\end{solution}

\begin{solution}{exo:b3:lithosphere-thermal-models:11}
El término $n = 3$ es $3104\times(8/\pi^2)\times\tfrac19e^{-9t/36.2} =
280\,e^{-9t/36.2}$ metros: \qty{218}{m} a \qty{1}{Myr}, \qty{23}{m} a
\qty{10}{Myr} y \qty{1.9}{m} a \qty{20}{Myr}. Desciende hasta \qty{10}{m} cuando
$e^{-9t/36.2} = 10/280$, en $t = (36.2/9)\ln28 = \qty{13.4}{Myr}$.
\end{solution}

\begin{solution}{exo:b3:lithosphere-thermal-models:12}
(a) Para $z \gg h_r$, $e^{-z/h_r} \to 0$. (b) $q_m = \qty{24}{mW.m^{-2}}$,
$A_0h_r^2/k = \num{16e-3}\times\num{e4}/2.5 = \qty{64}{K}$, y $z = (1300 -
64)\times2.5/\num{24e-3} = \qty{128.8}{km}$. (c) $q_m = 2.5\times1236/\num{250e3}
= \qty{12.4}{mW.m^{-2}}$: el \omterm{def:b3:lithosphere-thermal-models:reduced}{flujo de calor reducido} debe reducirse aproximadamente a la mitad, como ocurre bajo un
cratón cuya raíz gruesa y fría deja entrar poco calor desde abajo.
\end{solution}

\begin{solution}{pb:b3:lithosphere-thermal-models:1}
\textbf{1.} $\sqrt{\num{0.8e-6}\times\num{3.156e14}} = \qty{15889}{m}$ y
$\sqrt{\num{0.8e-6}\times\num{1.2624e15}} = \qty{31778}{m}$.

\textbf{2.} $\eta = 20000/31778 = 0.6294$ y $20000/63556 = 0.3147$:
$T = 1450\times0.6266 = \qty{909}{\celsius}$ a \qty{10}{Myr} y
$1450\times0.3437 = \qty{498}{\celsius}$ a \qty{40}{Myr}.

\textbf{3.} $z = 2\times0.7174\sqrt{\kappa t}$: \qty{22.8}{km} y
\qty{45.6}{km}. Alcanza \qty{60}{km} cuando $\sqrt{\kappa t} = 60/1.4348 =
\qty{41.82}{km}$, en $t = (\num{41.82e3})^2/\num{0.8e-6} = \qty{2.186e15}{s}$,
\qty{69.3}{Myr}.

\textbf{4.} $q = 3.1\times1450/(\sqrt\pi\times15889) = \qty{159.6}{mW.m^{-2}}$ a
\qty{10}{Myr} y \qty{79.8}{mW.m^{-2}} a \qty{40}{Myr}; $C_q = 159.6\sqrt{10}
= \qty{505}{mW.m^{-2}.Myr^{1/2}}$.

\textbf{5.} $u = \qty{7.92e-10}{m.s^{-1}}$, $\mathrm{Pe} = 7.92\times10^{-10}
\times\num{95e3}/\num{0.8e-6} = 94$. El calor viaja lateralmente con la placa a
cien veces más deprisa de lo que se difunde lateralmente, así que cada columna se enfría
solo por conducción vertical.

\textbf{6.} $(\num{95e3})^2/\num{0.8e-6} = \qty{1.13e16}{s}$, \qty{357}{Myr}.
A \qty{150}{Myr}, $\sqrt{\kappa t} = \qty{61.5}{km}$ y la capa enfriada, a
unas pocas veces esa cifra, llega mucho más abajo de \qty{95}{km}: un modelo de semiespacio para un fondo tan antiguo
supone que el frío ha alcanzado profundidades donde el manto convecta.

\textbf{7.} $(\rho_m - \rho_w)(d - d_r) = \rho_m\alpha\int(T_m - T)\,\mathrm dz$,
las masas de la columna de la dorsal y de la columna antigua siendo iguales por encima de una
profundidad de compensación. $T_m - T = T_m\operatorname{erfc}\eta$ cae por debajo de
\qty{0.5}{\percent} de $T_m$ para $\eta > 2$, por lo que la parte profunda no aporta nada.

\textbf{8.} $C = 2\times3300\times\num{3.1e-5}\times1450/2270\times
\sqrt{\num{0.8e-6}\times\num{3.156e13}/\pi} = 0.13069\times2834.8 =
\qty{370.5}{m.Myr^{-1/2}}$; $d = 2500 + 370.5\sqrt{10} = \qty{3672}{m}$ y
$2500 + 370.5\sqrt{40} = \qty{4843}{m}$.

\textbf{9.} El modelo da $2500 + 370.5\sqrt{42.5} = \qty{4915}{m}$,
\qty{440}{m} más profundo que la mediana. La mediana incluye la carga de los
sedimentos, los montes submarinos y las mesetas volcánicas, los abombamientos de los puntos calientes, y
las celdas en los flancos de zonas de fractura y márgenes, todos los cuales hacen que el fondo
más somero.

\textbf{10.} $C$ es proporcional a $T_m$: $370.5\times1350/1450 =
\qty{344.9}{m.Myr^{-1/2}}$, cercano al 350 ajustado.

\textbf{11.} El agua de mar que circula por la corteza joven fracturada transporta
parte del calor hacia el exterior a través de fuentes hidrotermales, entre las sondas.

\textbf{12.} $\tau = (\num{95e3})^2/(\pi^2\times\num{0.8e-6}) =
\qty{1.143e15}{s} = \qty{36.2}{Myr}$; la amplitud es $3300\times
\num{3.1e-5}\times1450\times\num{95e3}/4540 = \qty{3104}{m}$.

\textbf{13.} $e^{-40/36.2} = 0.33124$, $e^{-360/36.2}/9 = \num{5.3e-6}$ y el
término $n = 5$ es aún menor; el corchete es $1 - 0.81057\times0.33125 =
0.73150$, y $d = 2500 + 3104\times0.73150 = \qty{4771}{m}$.

\textbf{14.} Por el término $n = 3$, $3104\times0.81057\times\num{5.3e-6} =
\qty{0.013}{m}$: a \qty{40}{Myr}, $e^{-9t/\tau}$ es $10^{-4}$ de
$e^{-t/\tau}$, y solo queda el primer modo.

\textbf{15.} A \qty{40}{Myr}, $4843 - 4771 = \qty{73}{m}$; a
\qty{50}{Myr}, $2500 + 370.5\sqrt{50} = 5120$ frente a $2500 + 3104\times(1 -
0.81057\times0.25127) = 4972$: \qty{148}{m}. Interpolando, $40 + 10\times
(100 - 73)/(148 - 73) = \qty{43.6}{Myr}$, unos \qty{44}{Myr} (la serie completa
da 44.0).

\textbf{16.} $q = (3.1\times1450/\num{95e3})[1 + 2(0.33124 + 0.01204 +
0.00005)] = 47.32\times1.68666 = \qty{79.8}{mW.m^{-2}}$, la misma que la del
semiespacio: a \qty{40}{Myr} la base aún no se ha dejado sentir en la
superficie.

\textbf{17.} $2500 + 3104 = \qty{5604}{m}$ y $kT_m/L = \qty{47.3}{mW.m^{-2}}$.

\textbf{18.} $d(120) = 2500 + aL\bigl[1 - (8/\pi^2)e^{-120/\tau}\bigr]$, con
$a = \rho_m\alpha T_m/\bigl(2(\rho_m - \rho_w)\bigr) = \qty{32.673}{m/km}$ y
$\tau = L^2/(\pi^2\kappa)$.

\textbf{19.} $L = 3200/\bigl(32.673\,[1 - 0.81057e^{-120/\tau}]\bigr)$. Desde
$L = \qty{100}{km}$: $\tau = \qty{40.13}{Myr}$, corchete 0.95925, $L =
\qty{102.10}{km}$; luego $\tau = 41.83$, $L = 102.67$; $\tau = 42.30$, $L =
102.82$; $\tau = 42.43$, $L = 102.87$; se estabiliza en $L = \qty{102.9}{km}$.

\textbf{20.} $\tau = \qty{42.48}{Myr}$, $aL = \qty{3362}{m}$: $d(40) = 2500 +
3362\times(1 - 0.81057e^{-40/42.48}) = \qty{4799}{m}$, frente a \qty{4843}{m}
para el semiespacio y \qty{4800}{m} medido: la placa se ajusta a ambas edades.

\textbf{21.} $\tau = \qty{42.5}{Myr}$; profundidad asintótica $2500 + 3362 =
\qty{5862}{m}$.

\textbf{22.} $3362\times0.81057\,e^{-t/42.48} = 100$ da $t = 42.48\ln27.25 =
\qty{140}{Myr}$.

\textbf{23.} La placa que se ajusta al fondo antiguo del océano inventado es
\textbf{\qty[text-series-to-math=true]{102.9}{km}} de espesor.
\end{solution}
